\documentclass[11pt,a4paper]{article}
\usepackage[margin=1in,headheight=14pt]{geometry}
\usepackage{amsmath,amssymb}
\usepackage{xcolor,fancyhdr,enumitem}
\usepackage[hidelinks]{hyperref}
\pagestyle{fancy}
\fancyhf{}
\fancyhead[L]{\small\itshape Reinforcement Learning --- Fall 2026}
\fancyhead[R]{\small\itshape Homework 1: Introduction and Multi-armed Bandit}
\fancyfoot[C]{\thepage}
\renewcommand{\headrulewidth}{0.4pt}
\setlength{\parindent}{0pt}
\setlength{\parskip}{5pt}
\setlist[enumerate]{leftmargin=*,itemsep=9pt}
\newcommand{\answer}{\par\smallskip\textbf{Answer}\par}
\begin{document}
\begin{center}
{\Large\bfseries Reinforcement Learning --- Fall 2026\par}
\vspace{10pt}
{\LARGE\bfseries Homework 1: Introduction and Multi-armed Bandit\par}
\vspace{10pt}
{\large\bfseries\color{blue} Theoretical and Implementation Parts\par}
\vspace{10pt}
For any questions regarding this assignment, please contact\\
Amirparsa Bahrami at \texttt{amirparsa2828@gmail.com}.
\end{center}
\vspace{8pt}
{\large\bfseries\color{blue} Theoretical Part}
\section{Problem 1: Upper Confidence Bound (UCB) Action Selection}
The Upper Confidence Bound (UCB) algorithm selects actions according to:
\[
A_t=\arg\max_a\left[Q_t(a)+c\sqrt{\frac{\ln t}{N_t(a)}}\right]
\]
Explain intuitively:
\begin{enumerate}
\item What each term represents.
\answer
$A_t$ is the action selected at time $t$. $Q_t(a)$ estimates the mean reward of action $a$, favoring actions that have performed well (exploitation). $N_t(a)$ is the number of times $a$ has been selected before time $t$. The term $\sqrt{\ln t/N_t(a)}$ is an exploration bonus: actions with fewer observations receive a larger bonus because their values are more uncertain. The constant $c\geq0$ scales this bonus. The $\arg\max$ chooses the action with the largest sum of estimated value and exploration bonus. An untried action ($N_t(a)=0$) must be tried first, or equivalently assigned an infinite bonus.

\item Why the exploration bonus
\[
\sqrt{\frac{\ln t}{N_t(a)}}
\]
decreases with more samples.
\answer
For a fixed $t$, increasing $N_t(a)$ decreases the bonus proportionally to $1/\sqrt{N_t(a)}$. More observations reduce uncertainty, so less extra encouragement is needed to try that action. The numerator grows only logarithmically with time, allowing neglected actions to become attractive again. Thus, the decrease is with sample count at fixed time; the bonus need not decrease at every successive time step.

\item How the constant $c$ affects exploration.
\answer
A larger $c$ gives uncertainty more weight, encouraging more exploration and possibly sacrificing immediate reward to identify better actions. A smaller $c$ emphasizes current estimates and exploits more, but can take longer to correct inaccurate estimates. With $c=0$, the rule is greedy once untried actions have been handled.
\end{enumerate}

\newpage
\section{Problem 2: Gradient Bandit Algorithm}
In the gradient bandit algorithm, each action $a$ is assigned a preference $H_t(a)$, and the agent chooses actions using a softmax distribution:
\[
\pi_t(a)=\frac{e^{H_t(a)}}{\sum_b e^{H_t(b)}}
\]
\begin{enumerate}
\item Starting from the objective
\[
J(H)=\mathbb{E}[R_t],
\]
show that the gradient of the expected reward with respect to the preferences is:
\[
\nabla_{H_t(a)}J=\mathbb{E}[R_t(\mathbb{I}_{a=A_t}-\pi_t(a))]
\]
\answer
Let $q_*(b)=\mathbb{E}[R_t\mid A_t=b]$, which does not depend on the preferences. Then
\[
J(H_t)=\sum_b\pi_t(b)q_*(b).
\]
Differentiating the softmax gives
\[
\frac{\partial\pi_t(b)}{\partial H_t(a)}
=\pi_t(b)\bigl(\mathbb{I}_{a=b}-\pi_t(a)\bigr).
\]
Therefore,
\begin{align*}
\frac{\partial J}{\partial H_t(a)}
&=\sum_b q_*(b)\pi_t(b)\bigl(\mathbb{I}_{a=b}-\pi_t(a)\bigr)\\
&=\mathbb{E}\!\left[R_t\bigl(\mathbb{I}_{a=A_t}-\pi_t(a)\bigr)\right].
\end{align*}
The expectation is over the sampled action and its reward, for the current preferences.

\item From this result, derive the stochastic gradient ascent update rule used in practice.
\answer
Replace the expected gradient by its single-sample estimate and take an ascent step of size $\alpha>0$:
\[
H_{t+1}(a)=H_t(a)+\alpha R_t\bigl(\mathbb{I}_{a=A_t}-\pi_t(a)\bigr).
\]
An action-independent baseline $b_t$, fixed before sampling the current action and reward, can be subtracted without changing the expected gradient because
\[
\mathbb{E}\!\left[b_t(\mathbb{I}_{a=A_t}-\pi_t(a))\right]
=b_t\bigl(\pi_t(a)-\pi_t(a)\bigr)=0.
\]
Using the running average reward $\bar R_t$ available before this step yields
\[
\boxed{H_{t+1}(a)=H_t(a)+\alpha(R_t-\bar R_t)
\bigl(\mathbb{I}_{a=A_t}-\pi_t(a)\bigr),\quad\text{for every }a.}
\]
Afterward, the average may be updated as
$\bar R_{t+1}=\bar R_t+(R_t-\bar R_t)/t$, starting from $\bar R_1=0$.

\item Explain intuitively how this update changes the action preferences when $R_t$ is greater or less than the baseline $\bar R_t$.
\answer
For the selected action, the change is $\alpha(R_t-\bar R_t)(1-\pi_t(A_t))$; for every other action $a$, it is $-\alpha(R_t-\bar R_t)\pi_t(a)$. A reward above the baseline increases the selected action's preference and decreases the others, making the selected action more likely next time. A reward below the baseline reverses these changes, making it less likely. A reward equal to the baseline produces no change.
\end{enumerate}

\newpage
\section{Problem 3: Sample-Average and Constant Step-Size Updates}
For a fixed action, let $R_n$ denote the reward obtained on its $n$th selection, and let $Q_n$ denote its value estimate immediately before that selection.
\begin{enumerate}
\item Starting from the sample-average estimate, show that its incremental update (for $n\geq2$) is
\[
Q_{n+1}=Q_n+\frac{1}{n}(R_n-Q_n).
\]
\answer
Before the $n$th selection, the sample-average estimate is
\[
Q_n=\frac{1}{n-1}\sum_{i=1}^{n-1}R_i,\qquad n\geq2.
\]
Including the new reward gives
\begin{align*}
Q_{n+1}
&=\frac{1}{n}\sum_{i=1}^{n}R_i
=\frac{(n-1)Q_n+R_n}{n}\\
&=Q_n+\frac{1}{n}(R_n-Q_n).
\end{align*}
For the first observation, the step size is $1$, so $Q_2=R_1$ regardless of $Q_1$.

\item Now consider a constant step size $0<\alpha\leq1$:
\[
Q_{n+1}=Q_n+\alpha(R_n-Q_n).
\]
By expanding this recursion, show that
\[
Q_{n+1}=(1-\alpha)^nQ_1+\sum_{i=1}^{n}\alpha(1-\alpha)^{n-i}R_i.
\]
\answer
Write the update as $Q_{n+1}=(1-\alpha)Q_n+\alpha R_n$. Repeated substitution gives
\begin{align*}
Q_{n+1}
&=(1-\alpha)\bigl((1-\alpha)Q_{n-1}+\alpha R_{n-1}\bigr)+\alpha R_n\\
&=(1-\alpha)^2Q_{n-1}+\alpha(1-\alpha)R_{n-1}+\alpha R_n\\
&=\cdots\\
&=(1-\alpha)^nQ_1+\sum_{i=1}^{n}\alpha(1-\alpha)^{n-i}R_i.
\end{align*}
Each substitution multiplies older terms by $1-\alpha$. The weights sum to one:
\[
(1-\alpha)^n+\alpha\sum_{j=0}^{n-1}(1-\alpha)^j=1.
\]
When $\alpha=1$, the update simply gives $Q_{n+1}=R_n$, with all older terms having zero weight.

\item Using the weights in these two updates, explain which method is better suited to an environment where the mean rewards change over time. What trade-off arises when $\alpha$ is increased?
\answer
The sample average assigns every observed reward weight $1/n$, so old rewards remain as influential as recent ones and the step size shrinks toward zero. It therefore becomes slow to track changing means.

A constant step size is better suited to changing means: for $0<\alpha<1$, a reward $k$ selections old has weight $\alpha(1-\alpha)^k$, so outdated observations are progressively forgotten. Increasing $\alpha$ makes the estimate respond faster to changes but also more sensitive to random reward noise, producing greater fluctuations. A smaller $\alpha$ smooths the estimate but adapts more slowly; $\alpha=1$ uses only the latest reward.
\end{enumerate}

\newpage
\section{Problem 4: Optimistic Initial Action Values}
Consider a deterministic, stationary three-armed bandit with rewards $r(1)=1$, $r(2)=2$, and $r(3)=3$. The agent uses purely greedy action selection ($\varepsilon=0$), updates selected actions using their sample-average rewards, and breaks ties by choosing the smallest action index.
\begin{enumerate}
\item Suppose the initial estimates are optimistic: $Q_1(a)=5$ for every action $a$. Determine the first four selected actions and the value estimates after the third selection.
\answer
The first observation of an action replaces its initial estimate by its reward. Starting from $(5,5,5)$:
\begin{center}
\begin{tabular}{c c c l}
\hline
Selection & Chosen action & Estimates after selection & Reason for choice\\
\hline
1 & 1 & $(1,5,5)$ & All tied; smallest index is 1.\\
2 & 2 & $(1,2,5)$ & Actions 2 and 3 tied at 5.\\
3 & 3 & $(1,2,3)$ & Action 3 has largest estimate.\\
4 & 3 & $(1,2,3)$ & Action 3 has largest estimate.\\
\hline
\end{tabular}
\end{center}
Thus, the first four actions are $\boxed{1,2,3,3}$, and immediately after the third selection the estimates are $\boxed{Q_4=(1,2,3)}$.

\item Repeat the first-four-actions calculation when $Q_1(a)=0$ for every action. Explain why the behavior differs.
\answer
Starting from $(0,0,0)$, the tie-breaking rule selects action 1 first. Its estimate becomes 1, while the other two remain 0:
\begin{center}
\begin{tabular}{c c c}
\hline
Selection & Chosen action & Estimates after selection\\
\hline
1 & 1 & $(1,0,0)$\\
2 & 1 & $(1,0,0)$\\
3 & 1 & $(1,0,0)$\\
4 & 1 & $(1,0,0)$\\
\hline
\end{tabular}
\end{center}
The first four actions are $\boxed{1,1,1,1}$. Every reward from action 1 is 1, so its sample average stays 1 and the greedy agent never tries actions 2 or 3. With optimistic initialization, untried actions instead retain estimates above the observed rewards, which drives the agent to try them.

\item Explain how optimistic initialization can encourage exploration even with $\varepsilon=0$. Why might this approach be insufficient when the true action values change over time?
\answer
Optimistic estimates make untried or insufficiently tried actions look attractive. Selecting an action and observing a reward below its optimistic estimate lowers its estimate, allowing other actions to become greedy choices. This creates exploration without random action selection.

This incentive is temporary: once the initial optimism has been replaced by observations, a greedy agent has no built-in reason to revisit actions with low estimates. If a neglected action later improves, the agent may never discover the change. Sample-average estimates also adapt slowly after many observations. Continued exploration, together with an update that can track changes, is therefore needed for reliable adaptation.
\end{enumerate}
\end{document}
